% 原书物理页=印刷页45--48；1.13 全节.
\section{$\mathfrak{sl}_2(\mathbb R)$ 及其化身的不可约表示}\label{sec:lie-1-13}

我们将在本节对 $\mathfrak{sl}_2(\mathbb R)$ 的不可约复表示进行分类，进而导出 $\mathfrak{sl}_2(\mathbb C)$ 以及群 $\mathrm{SL}_2(\mathbb R)$、$\mathrm{SU}(2)$、$\mathrm{SL}_2(\mathbb C)$ 与 $\mathrm{SO}_3(\mathbb R)$ 的不可约表示的分类.\footnote{译注：节标题中的\newterm{lie-avatar}{化身}{avatar}（avatars），指 $\mathfrak{sl}_2(\mathbb R)$ 的各种实、复实现，如 $\mathfrak{sl}_2(\mathbb C)$、$\mathfrak{su}(2)$、$\mathfrak{so}_3(\mathbb R)$ 等，见下文 \ref{np:lie-1-13-2} 与 \ref{np:lie-1-13-3} 段.}
我们已经看到，所有这些对象的表示都是\emph{完全可约}的，因此在相差一个同构的意义下，上述分类描述了这些对象的全部表示.

为分类 $\mathfrak{sl}_2(\mathbb R)$ 的不可约表示而使用的技巧与中间结果，对我们随后研究``半单''Lie 代数将是根本性的.

\begin{numpara}[$\mathfrak{g}=\mathfrak{sl}_2(\mathbb R)$ 的不可约表示]\label{np:lie-1-13-1}
回顾 $\dim_{\mathbb R}\bigl(\mathfrak{sl}_2(\mathbb R)\bigr)=3$. 我们将使用下面这组基 $(E,H,F)$：
\[
E=\begin{pmatrix}0&1\\0&0\end{pmatrix},\qquad
H=\begin{pmatrix}1&0\\0&-1\end{pmatrix},\qquad
F=\begin{pmatrix}0&0\\1&0\end{pmatrix}.
\]
$\mathfrak{g}:=\mathfrak{sl}_2(\mathbb R)$ 的 Lie 括号在这组基下由下列公式确定：
\[
[H,E]=2E,\qquad [H,F]=-2F,\qquad [E,F]=H.
\]
\par
若 $(V,r)$ 是 $\mathfrak{g}=\mathfrak{sl}_2(\mathbb R)$ 的一个表示且 $\lambda\in\mathbb C$，我们记
\[
V_\lambda:=\operatorname{Ker}\bigl(r(H)-\lambda\,\mathrm{id}_V\bigr).
\]
若 $V_\lambda\neq 0$，它就是 $r(H)$ 的属于特征值 $\lambda$ 的特征空间. 此时我们称 $\lambda$ 是 $H$ 在 $V$ 中的一个\newterm{lie-weight}{权}{weight}.
\end{numpara}

\begin{lemma*}
设 $(V,r)$ 是 $\mathfrak{g}$ 的一个复表示，$\lambda\in\mathbb C$ 是 $H$ 在 $V$ 中的一个权. 则 $r(E)V_\lambda\subset V_{\lambda+2}$ 且 $r(F)V_\lambda\subset V_{\lambda-2}$.
\end{lemma*}

\begin{proof}
设 $v\in V_\lambda$. 我们有 $HEv=EHv+[H,E]v=\lambda Ev+2Ev=(\lambda+2)Ev$. 同理 $HFv=(\lambda-2)Fv$.
\end{proof}

\begin{proposition*}
对任意整数 $m\in\mathbb N$，存在唯一的（在相差一个同构的意义下）$\mathfrak{g}$ 的 $m+1$ 维不可约复表示 $(V_m,r_m)$. 此外，$H$ 在 $V_m$ 中的权为 $\{-m,-m+2,\cdots,m-2,m\}$.
\end{proposition*}

\begin{proof}
i) 先证唯一性. 设 $V$ 是 $\mathfrak{g}$ 的一个不可约复表示，$\lambda_0$ 是一个实部最大的权，$v_0\in V_{\lambda_0}$ 非零. 根据上一引理，此时有 $EV_{\lambda_0}=0$.
令 $v_k:=\frac{1}{k!}F^k v_0$，并设 $W=\operatorname{Vect}_{\mathbb C}\{v_k,\,k\in\mathbb N\}$ 为由诸 $v_k$ 张成的空间. 由构造，该空间在 $F$ 下稳定；又由于诸 $v_k$ 都是 $H$ 的特征向量（上一引理），它也在 $H$ 下稳定. 现在证明它在 $E$ 下稳定. 前面已见到 $Ev_0=0$. 对 $k>1$，我们有
\[
Ev_k=\frac{1}{k}EFv_{k-1}=\frac{1}{k}\bigl(FEv_{k-1}+[E,F]v_{k-1}\bigr)=\frac{1}{k}\bigl(FEv_{k-1}+(\lambda_0-2k+2)v_{k-1}\bigr).
\]
这使我们能够用归纳法证明 $Ev_k=(\lambda_0-k+1)v_{k-1}$.
由此得出 $W$ 是 $V$ 的一个非零子表示，故由不可约性知 $W=V$. 非零的诸 $v_k$ 线性无关，因为它们对应于 $H$ 的两两不同的特征值，从而（令 $m=\dim_{\mathbb C}(V)-1$）$v_0,\cdots,v_m$ 是 $V$ 的一组基.
为确定 $\lambda_0$，我们注意到：由于 $H=[E,F]$，必须有 $\operatorname{tr}(r(H))=0$. 然而 $\operatorname{tr}(r(H))=\lambda_0+(\lambda_0-2)+\cdots+(\lambda_0-2m)=(m+1)\lambda_0-m(m+1)$. 于是 $\lambda_0=m$，权最终正是所宣告的那些.
ii) 现在需要验证存在性. 为此，只需验证 $\mathrm{M}_{m+1}(\mathbb C)$ 中下列矩阵
\[
H_m:=\sum_{i=1}^{m+1}(m-2i+2)E_{i,i},\qquad
E_m:=\sum_{i=1}^{m}(m-i+1)E_{i,i+1},\qquad\text{以及}\quad
F_m:=\sum_{i=1}^{m}i\,E_{i+1,i}
\]
满足关系 $[H_m,E_m]=2E_m$、$[H_m,F_m]=-2F_m$ 与 $[E_m,F_m]=H_m$. 此后只需令 $r_m(H):=H_m$、$r_m(E):=E_m$ 且 $r_m(F):=F_m$ 即可.
\end{proof}

\begin{remark*}
在小维度时，可以认出一些熟悉的表示：
\begin{itemize}[label=--]
\item $(V_0,r_0)$ 是 $\mathfrak{g}$ 的单位表示.
\item $(V_1,r_1)$ 是由包含关系 $\mathfrak{g}\subset\mathfrak{gl}_2(\mathbb C)$ 给出的``\newterm{lie-standard-representation}{标准表示}{standard representation}''.
\item $(V_2,r_2)$ 是（复化的）伴随表示 $\mathfrak{g}\longrightarrow\mathfrak{gl}(\mathfrak{g}_{\mathbb C})\simeq\mathfrak{gl}_3(\mathbb C)$.
\end{itemize}
一般地，可以证明 $(V_m,r_m)\simeq\operatorname{Sym}^m(V_1,r_1)$，即 $(V_1,r_1)$ 的 $m$ 次\newterm{lie-symmetric-power}{对称幂}{symmetric power}，其底空间为 $\operatorname{Sym}^m(V_1)=(V_1\otimes V_1\otimes\cdots\otimes V_1)^{\mathfrak{S}_m}$，而 $\mathfrak{g}$ 的作用是张量积表示的商.
\end{remark*}

\begin{remark*}
表示 $(V_m,r_m)$ 称为\newterm{lie-highest-weight}{最高权}{highest weight}为 $m$ 的不可约表示. 注意到，证明中用到的基 $E,H,F$ 的性质仅有公式 $[H,E]=2E$、$[H,F]=2F$\footnote{译注：按 $\mathfrak{sl}_2(\mathbb R)$ 的括号关系应为 $[H,F]=-2F$，原书如此.} 与 $[E,F]=H$. 在 $\mathrm{GL}_2(\mathbb R)$ 的作用下，任何满足这些性质的另外的基 $(E',H',F')$ 都共轭于 $(E,H,F)$. 实际上，把上述证明 ii) 中的公式应用到标准表示上，我们就得到 $\mathbb R^2$ 的一组基 $v_0,v_1$，使得 $(E',H',F')$ 在这组基下的矩阵恰好是 $E,H,F$.

此外，对任意实数 $\lambda>0$，都可以构造最高权为 $\lambda$ 的不可约表示，但其维数为无穷……
\end{remark*}

\begin{exercise*}
设 $(V,r)$ 是 $\mathfrak{sl}_2(\mathbb R)$ 的一个表示.
\begin{enumerate}[label=\roman*)]
\item 证明 $V$ 是它的各权子空间的和，它的权都是整数，并且\newterm{lie-multiplicity}{重数}{multiplicity}函数 $\lambda\mapsto\operatorname{mult}(\lambda,V):=\dim(V_\lambda)$ 在限制到 $2\mathbb Z$ 或 $2\mathbb Z+1$ 上时，是绝对值 $|\lambda|$ 的递减函数.
\item 证明 $V\simeq\bigoplus_{m\in\mathbb N}V_m^{\oplus e_m}$，其中 $e_m=\operatorname{mult}(m,V)-\operatorname{mult}(m+2,V)$.
\end{enumerate}
\end{exercise*}

\begin{numpara}[相关代数的不可约表示]\label{np:lie-1-13-2}
\begin{enumerate}[label=\roman*)]
\item $\mathbb C$-代数 $\mathfrak{sl}_2(\mathbb C)$. 根据到目前为止的结果，$\mathfrak{sl}_2(\mathbb C)$ 的``全纯''不可约表示按其最高权 $m\in\mathbb N$ 分类，并由与 $\mathfrak{sl}_2(\mathbb R)$ 相同的公式（$r_m(H)=H_m$，等等）经 $\mathbb C$-线性扩张给出. ``反全纯''的不可约表示同理给出，但经 $\mathbb C$-半线性扩张；而任何不可约表示都形如 $(V_m,r_m)\otimes(V_{m'},r_{m'})$.
\item $\mathbb R$-代数 $\mathfrak{su}(2)$. 任何不可约表示都可由 $\mathfrak{sl}_2(\mathbb C)$ 的一个全纯不可约表示限制得到，因而形如 $(V_m,r_m)$. 注意，$\mathfrak{su}(2)$ 有一组便于使用的基
\[
X_1=\begin{pmatrix}0&1\\-1&0\end{pmatrix},\qquad
X_2=\begin{pmatrix}0&i\\i&0\end{pmatrix},\qquad
X_3=\begin{pmatrix}i&0\\0&-i\end{pmatrix},
\]
其 Lie 括号由 $[X_1,X_2]=2X_3$、$[X_2,X_3]=2X_1$ 与 $[X_3,X_1]=2X_2$ 给出. 读者可以显式计算这些元素在 $V_m$ 上的作用. 特别地，$X_3$ 的特征值为 $-im,\,-i(-m+2),\,\cdots,\,im$.
\item $\mathbb R$-代数 $\mathfrak{so}_3(\mathbb R)$. 我们知道它与 $\mathfrak{su}(2)$ 同构，故其不可约表示仍由整数 $m$ 参数化. $\mathfrak{so}_3(\mathbb R)$ 有一组便于使用的基
\[
J_x=\begin{pmatrix}0&0&0\\0&0&-1\\0&1&0\end{pmatrix},\qquad
J_y=\begin{pmatrix}0&0&1\\0&0&0\\1&0&0\end{pmatrix},\qquad
J_z=\begin{pmatrix}0&-1&0\\1&0&0\\0&0&0\end{pmatrix},
\]
其 Lie 括号由 $[J_x,J_y]=J_z$、$[J_y,J_z]=J_x$ 与 $[J_z,J_x]=J_y$ 给出.\footnote{译注：按括号关系 $[J_x,J_y]=J_z$（并对照同态 $X_2\mapsto 2J_y$）验证，$J_y$ 在 $(3,1)$ 位置应为 $-1$；原书如此.} 由 $\mathfrak{su}(2)$ 的伴随表示给出的同构 $\mathfrak{su}(2)\stackrel{\sim}{\longrightarrow}\mathfrak{so}_3(\mathbb R)$ 把 $X_1$ 映到 $2J_x$、$X_2$ 映到 $2J_y$、$X_3$ 映到 $2J_z$. 特别地，$J_z$ 在 $V_m$ 中的特征值为 $-\frac{i}{2}m,\,\frac{i}{2}(-m+2),\,\cdots,\,\frac{i}{2}m$. 物理学家把 $m/2$ 称为表示 $V_m$ 的``\newterm{lie-spin}{自旋}{spin}''.
\end{enumerate}
\end{numpara}

\begin{numpara}[相关 Lie 群的表示]\label{np:lie-1-13-3}
理论告诉我们，$\mathfrak{sl}_2(\mathbb R)$、$\mathfrak{su}(2)$ 与 $\mathfrak{sl}_2(\mathbb C)$ 的表示 $(V_m,r_m)$ 分别可提升到群 $\mathrm{SL}_2(\mathbb R)$、$\mathrm{SU}(2)$ 与 $\mathrm{SL}_2(\mathbb C)$ 上，并给出全部的不可约表示（对 $\mathrm{SL}_2(\mathbb C)$ 而言是全纯的）. 然而，这并没有给我们这些提升的显式公式. 因此，为这些提升寻找一个自然的模型是合情合理的. 下面便是这样一个模型.

设 $\mathcal{P}_m:=\{f(x,y)\in\mathbb C[x,y],\ f \text{ 为次数 } m \text{ 的齐次多项式}\}$. 我们让 $\mathrm{SL}_2(\mathbb C)$ 按公式
\[
\Bigl[\begin{pmatrix}a&b\\c&d\end{pmatrix}f\Bigr](x,y):=f(ax+cy,\; bx+dy)
\]
作用. 这样便得到 $\mathrm{SL}_2(\mathbb C)$ 的一个全纯表示，我们记之为 $(\mathcal{P}_m,\rho_m)$.
\end{numpara}

\begin{lemma*}
$(\mathcal{P}_m,d\rho_m)\simeq(V_m,r_m)$.
\end{lemma*}

\begin{proof}
我们计算（与 \ref{np:lie-1-8-1} 段的诸例子相比较）
\begin{align*}
[d\rho_m(H)f](x,y)&=\frac{d}{dt}\Big|_{t=0}\bigl(t\mapsto\rho_m(\exp(tH))f(x,y)\bigr)\\
&=\frac{d}{dt}\Big|_{t=0}\bigl(t\mapsto f(\exp(t)x,\exp(-t)y)\bigr)\\
&=\Bigl(x\frac{\partial}{\partial x}-y\frac{\partial}{\partial y}\Bigr)(f)(x,y).
\end{align*}
于是令 $f_j:=x^{m-j}y^j$. 族 $f_0,\cdots,f_m$ 是 $\mathcal{P}_m$ 的一组基，并且我们发现 $d\rho_m(H)f_j=(m-2j)f_j$. 因此 $\mathcal{P}_m$ 的权是 $-m,\,-m+2,\,\cdots,\,m-2,\,m$，且重数均为 1（特征空间的维数为 1）. 这足以断言 $(\mathcal{P}_m,d\rho_m)=(V_m,r_m)$. 实际上，我们知道 $(\mathcal{P}_m,d\rho_m)$ 同构于不可约表示的一个直和 $\bigoplus_n V_n^{\oplus e_n}$. 由于 $\dim\mathcal{P}_m=m+1$，可见对 $n>m$ 有 $e_n=0$. 而既然 $m$ 是 $\mathcal{P}_m$ 的一个权，必有 $e_m\geqslant 1$，再由维数的原因便得 $e_m=1$ 且对 $n\neq m$ 有 $e_n=0$.（也可以应用上面习题的 ii).）
\end{proof}

通过限制，$(\mathcal{P}_m,\rho_m)$ 给出 $\mathrm{SL}_2(\mathbb R)$ 与 $\mathrm{SU}(2)$ 的不可约表示. 为得到 $\mathrm{SO}_3(\mathbb R)$ 的那些表示，我们利用二叶覆盖 $\mathrm{SU}(2)\longrightarrow\mathrm{SO}_3(\mathbb R)$，其核为 $\{\pm I_2\}$. 它表明，$\mathrm{SO}_3(\mathbb R)$ 的表示可等同于 $\mathrm{SU}(2)$ 的满足 $\rho(-I_2)=\mathrm{id}_V$ 的那些表示 $(V,\rho)$. 于是可以验证（习题）：$(\mathcal{P}_m,\rho_m)$ ``下降''到 $\mathrm{SO}_3(\mathbb R)$ 上当且仅当 $m$ 是偶数.

\begin{corollary*}
$\mathrm{SO}_3(\mathbb R)$ 的不可约复表示的维数都是奇数，并且由该维数唯一确定.
\end{corollary*}

特别地，我们看到 $\mathfrak{so}_3(\mathbb R)$ 的表示并非全体都能积分成 $\mathrm{SO}_3(\mathbb R)$ 的表示.

\begin{numpara}[$\mathrm{SO}_3(\mathbb R)$ 在球面 $S^2$ 上函数的作用]\label{np:lie-1-13-4}
前面已经解释过，$\mathrm{SO}_3(\mathbb R)$ 在球面 $S^2$ 上的作用如何诱导出 $C^\infty(S^2)$ 上的一个作用（这里取复值函数）. 我们记 $\mathcal{P}^{\mathrm{sph}}_{\leqslant n}$ 为 $C^\infty(S^2)$ 的由下述函数组成的子空间：它们是 $x,y,z\in\mathbb R^3$ 的次数 $\leqslant n$ 的复系数多项式限制到球面上所得的函数.
\end{numpara}

\begin{proposition}\label{prop:lie-1-13-5}
表示 $\mathcal{P}^{\mathrm{sph}}_{\leqslant n}$ 同构于 $V_0\oplus V_2\oplus\cdots\oplus V_{2n}$，其中 $V_{2m}$ 表示 $\mathrm{SO}_3(\mathbb R)$ 的唯一的 $2m+1$ 维不可约表示.
\end{proposition}

\begin{proof}
令 $u=x+iy$，$v=x-iy$. 等式 $1-z^2=x^2+y^2=uv$ 表明，$\mathcal{P}^{\mathrm{sph}}_{\leqslant n}$ 由函数 $z^p$、$z^pu^k$ 与 $z^pv^k$（$p+k\leqslant n$）生成.

$\mathfrak{so}_3(\mathbb R)$ 在 $C^\infty(S^2)$ 上的导出作用已在 \ref{np:lie-1-8-1} 段的第二个例子中算出. 其中 $J_z$ 的作用由 $x\frac{\partial}{\partial y}-y\frac{\partial}{\partial x}=iu\frac{\partial}{\partial u}-iv\frac{\partial}{\partial v}$ 给出. 于是我们得到
\[
J_z(z^pu^k)=ik\,(z^pu^k)\qquad\text{以及}\qquad J_z(z^pv^k)=-ik\,(z^pv^k).
\]
因此，$J_z$ 在 $\mathcal{P}^{\mathrm{sph}}_{\leqslant n}$ 中的特征值形如 $-in,\,i(-n+1),\,\cdots,\,i(n-1),\,in$，从而能够出现在 $\mathcal{P}^{\mathrm{sph}}_{\leqslant n}$ 中的 $\mathfrak{sl}_3(\mathbb R)$\footnote{译注：按上下文应为 $\mathfrak{so}_3(\mathbb R)$，原书如此.} 的不可约表示只有诸 $V_{2m}$（$m\leqslant n$）. 现在注意，对固定的 $k$，函数 $u^k,\,zu^k,\,\cdots,\,z^{n-k}u^k$ 线性无关. 这表明 $J_z$ 的特征值 $ik$（从而权 $2k$）的重数是 $n-k$. 利用上面的最后一个习题，即得所宣告的分解.
\end{proof}
